\documentclass[12pt]{article}

\title{Week Two Lecture Notes}
\author{
	Evan M Drake \\
	\texttt{drakeev@butte.edu}
}


\usepackage{snow-preamble}

%%%%%%%%%%%%%%%%%%%%%%%%%%%%%%%%%%%%
\begin{document}
	\maketitle
	\tableofcontents

	\newpage

	%% body start

	\section{Data Types in MATLAB}
		In the below we can see a graphic depicting the type tree of MATLAB.
		While giving you a good picture of how the types are related and what exactly they mean,
			this can look complicated.
		The idea to understand is that it is generally the case that specific data types such as numerics,
			often give rise vase to bodies of data types such as floating-point (real numbers),
			integers,
			complex numbers,
			which creates type hierarchies within languages that can be exploited when optimizing. \\

		\includegraphics[width=\textwidth]{DT}

	\section{Common operators}
		Arithmetic is an important capability of MATLAB.
		As such the language contains all common arithmetic operators:
		\begin{itemize}
			\item Association: ( )
			\item Unary: positive +, negative -
			\item Arithmetic: \verb| +, -, *, /, ^ log, \% |
			\item Relational: \verb| <, >, <=, >=, ==, ~= |
			\item Logical (Bitwise): \verb| &, &&, | |, ||
		\end{itemize}

	\subsection{Operator Precedence}
		\begin{enumerate}
			\item ( )
			\item \verb| ^ .^ ’ .’ (pure transpose)|
			\item \verb| + (unary plus) - (unary minus) ~ (NOT)|
			\item \verb| * / \ .* ./ . |
			\item \verb| + (addition) - (subtraction)|
			\item \verb| : |
			\item \verb| > < >= <= == ~=|
			\item \verb| & (AND)|
			\item | (OR)
			\item \verb| && (scalar AND) |
			\item || (scalar OR)
			\item \verb| = (assignment) |
		\end{enumerate}

	\section{Variables}
		A variable is a \emph{value} bound to some arbitrary \emph{name}.
		It is the case that the name must be some \emph{word} (bound by the English alphabet and Arabic numerals in our case).
		Recall a we use word here as in set theory in which a word is some combination of values from a set.
		For example "and" is a word in English b/c $[a, n, d]$ are all from the set of English letters.
		Note that in MATLAB the names must start with a letter but can then contain arbitrary numbers.
		\subsection{Values / Expressions}
			Here a value is represented by any expression.
			An expression can be thought of as a kind of solvable,
				sequence of operations (equation),
				in discrete space\footnote{we must always be in discrete space on a computer}.
			So an example of an expression is $f(x) = x+5$ where $x+5$ is our expression.
			This represents a solvable,
				sequence of operations that is seemingly in discrete space.
		\subsection{The assignment operator}
			The assignment operator $=$ is how we will create variables in MATLAB.
			Note that in pseudocode they are often created using the binding operator $\leftarrow$.
	
	\section{Collections}
		A collection can be thought of as simply a homogeneous "bag of stuff".
		Where everything in the bag is of the same type.
		We can think of a bag of marbles,
			each "thing" in the bag is a glass bead but maybe they are of varying sizes.
		We can extend the idea of collections to many different structures in type theory.
		But it is important to understand collections for pseudo code purposes as they can assist in optimization.

		\subsection{Array}
			An array is simply an ordered collection of finite, predetermined size.
			In particular this means that we must know the size of an array before we create it.
			An array has order and therefore each item in it has an index with which we can find it.

		\subsection{Vector}
			A vector is a bit different from an array in that it has no predetermined size.
			An array has dynamic size,
			in particular the size of a vector can be augmented on the fly for convenience.
			Vectors also have a concept of \emph{direction},
				this concept should be familiar from linear algebra.
			We can create row vectors and/or column vectors.
			Thus we can also create matrices.

		\subsection{Matrix}
			The bases of linear algebra is the matrix and thus the we will define it here.
			A matrix is such that it contains both row and column vectors,
				giving rise to the common visualization of a kind of box of squares.
			Note that a matrix can be given a flattened visualization that is simply a vector of vectors.

		\subsection{Collection operations}
			\begin{itemize}
				\item Array Arithmetic: \verb| +, -, *, / | are mapped over the array space
				\item Vector arithmetic: inherits array arithmetic with extra rules for direction
				\item Matrix math: \emph{transpose}, \emph{multiplication} are the big ones here
			\end{itemize}

			Vector math can only be done between two vectors of the same length and direction.
			We wont worry about Cross multiplying as the calculation is easy to look up but difficult to memorize in our context.

			\begin{align*}
				\vec{u} + \vec{v} &\defeq \lbrack u_1+v_1,u_2+v_2...\rbrack \\
				\vec{u} - \vec{v} &\defeq \lbrack u_1-v_1,u_2-v_2...\rbrack \\
				\vec{u}\alpha &\defeq \lbrack u_1\alpha,u_2\alpha...\rbrack \\
				\vec{u} * \vec{v} &\defeq \sum_i^n u_iv_i \text{s.t} \\
				\text{\emph{Euclidean length} }\vec{u} &\defeq \sqrt{\sum_i^n u_i^2} \\
				\text{\emph{where} } n &\implies (|v| = |u|)
			\end{align*}

			The above are some of the common vector operations we will see in homework and labs.
			The most common Matrix operation in linear algebra is matrix multiplication as such we will deal with that often.
			Now I will urge you to revisit some linear algebra,
				I keep a reference with me and will probably have it in class often \cite{savov}.

			\begin{align*}
				A = &\mathbb{R}^{m x l} \\
				B = &\mathbb{R}^{l x n} \\
				|A * B| &\defeq \mathbb{R}^{m x n} \\
				A * B &\defeq A_R*B_C + A_{R+1}B_{C+1} + ...
			\end{align*}

			This will make much more sense after some muscle memory is used in the homework and assignments.

			\section{Linear Algebra}
				
				\subsection{Matrix Algebra}

				\subsection{Reduced Row echelon form}

				\subsection{Guass-Jordan Elimination}


	%% body finish

	\cite{*}

	\newpage % bib
	\backmatter
	\addcontentsline{toc}{section}{References}
	\printbibliography

	\newpage % index
	\printindex

	\newpage % appendices
	\appendix
		\section{Helpful Links}
		\href{https://www.mathworks.com/content/dam/mathworks/fact-sheet/matlab-basic-functions-reference.pdf}{List of common commands}

\end{document}

